\section{Time frequency transform with frequency focus}\label{sec_frequency_focus}
\label{sec:freqfocus}
It is also interesting to introduce frequency-dependent focus, in addition to time-dependent focus. We first stress that the construction of Section~\ref{sec_time_focus} may easily be transposed to that context. Indeed, given the symmetry property of the STFT provided by the Plancherel formula $\langle f,h_{t,\omega}\rangle = \langle \hat f,\widehat{h_{t,\omega}}\rangle = \langle\hat f,\hat h_{\omega,-t}\rangle$, "time-focus" may be applied to the STFT of the Fourier transform $\hat f$ of a signal $f$, resulting in frequency focus. Results similar to the ones described above can be obtained using the very same techniques, we won't address this adaptation here.

We shall rather address the introduction of frequency focus into another transform, which uses scale variables in place of frequency variables, namely wavelet and/or constant-Q transforms. These closely related transforms are based upon time-frequency atoms which have the constant-Q property. The Q factor is usually defined as the ratio of the central frequency $\xi$ of the atom by its spectral bandwidth $\delta\xi$ (both quantities will be properly defined below).

\subsection{Continuous constant-Q and wavelet transforms}
\subsubsection{Transforms on $L^2(\R)$}
The constant $Q$ transform was introduced in a discrete context~\cite{Brown1991calculation} and revisited more recently~\cite{velasco2011constructing,holighaus2012framework}. We provide below a slightly more general version adapted to the continuous setting.

The time-frequency atoms are built from a reference waveform $h\in L^2(\R)$, which will be assumed continuous and compactly supported in the Fourier domain.
Following the definition from~\cite{holighaus2012framework,velasco2011constructing}, time-frequency atoms $h_{t,\u}$ are generated as rescaled and shifted copies of $h$, which is implemented in the continuous setting as
\begin{equation}
	\forall x,t,\u\in\R\ ,\quad \ h_{t,u}(x) := \sqrt{\gamma(u)}e^{2i\pi\gamma(u)x}h(\gamma(u)x-t)\ .
\end{equation}
Here, $\gamma$ is a $C^1$ diffeomorphism such that $\lim_{u\to-\infty}\gamma(u)=0$ and $\lim_{u\to+\infty}\gamma(u)=+\infty$. In~\cite{holighaus2012framework,velasco2011constructing}, $\gamma$ was given an exponential form, we consider here a slightly more general such scale function. Time-frequency atoms $h_{t,u}$ are normalized so that $\|h_{t,u}\|_{L^2}=\|h\|_{L^2}$ for all $t,u$. The corresponding constant-Q transform maps every $f\in L^2(\R)$ to the function $Lf$ defined by
\begin{equation}
L_f(t,u)=\langle f,h_{t,u}\rangle_{L^2}\ .
\end{equation}
Since $f,h\in L^2(\R)$, $L_f(t,u)$ is well-defined for all $t,u\in\R$. Under suitable assumptions on $h$, $L$ also establishes an isometry between $L^2(\R)$ and $L^2(\R^2,d\mu)$, where the measure $\mu$ is defined by
\begin{equation}
\label{fo:CQT.measure}
d\mu(t,\u) := \tfrac{\gamma'(\u)}{\gamma(\u)}d\u dt \ .
\end{equation}
\begin{prop}\label{prop_norm_relation_operator_Q_transform}
Let $\gamma$ be a $C^1$ diffeomorphism such that $\lim_{\u\to -\infty}\gamma(\u)=0$ and $\lim_{\u\to +\infty}\gamma(\u)=+\infty$, let $h\in L^2(\R)$ satisfying the admissibility condition
\begin{equation}
\label{fo:CQT.admissibility}
    0<c_h := \int_{-1}^{+\infty}\frac{|\hat{h}(y)|^2}{y+1}dy = \int_{-\infty}^{-1}\frac{|\hat{h}(y)|^2}{-y-1}dy<\infty\ .
\end{equation}
Then for any $f\in L^2(\R)$ we have
\begin{equation}
\label{fo:CQT_isometry}
    \|Lf\|^2_{L^2(\R^2,d\mu)} = c_h \|f\|_{L^2(\R)}^2\ .
\end{equation}
\end{prop}
\begin{proof}[Sketch of the proof]
Let $h\in L^2(\R)$ satisfying the admissibility condition~\eqref{fo:CQT.admissibility}. Assume $f\in C_c(\R)$. Then $f\in L^1(\R)$, and by Young's convolution inequality $Lf(\cdot,u)\in L^2(\R)$ for all $u\in\R$. Introducing the auxiliary function $F(\xi,u)=\hat f(\xi) \overline{\hat h}(\xi/\gamma(u),u)$, we have
\[
\int_{\R}\! |Lf(t,u)|^2 dt = \int_{\R^3}\!\! F(\xi,u)\overline{F}(\xi',u)e^{2i\pi (\xi-\xi')t/\gamma(u)}\,d\xi d\xi' dt=\int_{\R}\! \left|\widecheck{F}_1(t,u)\right|^2 dt = \left\|F(\cdot,u)\right\|^2_{L^2}\ ,
%\int_\R \left|F(\xi,u)\right|^2\,du\ ,
\]
where we have used twice Plancherel's formula, and denoted by $\widecheck{F}_1=\mathcal{F}_1^{-1}F$ the inverse Fourier transform of $F$ with respect to its first variable.

Let $\varepsilon>0$, $I_\varepsilon = [\gamma^{-1}(1/\varepsilon),\gamma^{-1}(\varepsilon)]$. Focusing on positive frequencies first,  consider the (convergent) integral
\begin{eqnarray*}
\int_{I_\varepsilon} \|F(\cdot,u)\|_{L^2(\R_+)}^2 \frac{\gamma'(u)}{\gamma(u} du &=& \int_{I_\varepsilon\times\R_+} |\hat f(\xi)|^2\left|\hat\psi\left(\frac{\xi}{\gamma(u)}-1\right)\right|^2d\xi\frac{\gamma'(u)}{\gamma(u)}du\\
&=& \int_{\R_+} |\hat f(\xi)|^2\int_{-1+\xi\varepsilon}^{-1+\xi/\varepsilon} |\hat h(y)|^2\frac{dy}{y+1}\ .
\end{eqnarray*}
The inner integral is bounded by the admissibility constant $c_h$, the dominated convergence theorem then yields
\[
\int_\R \|F(\cdot,u)\|_{L^2(\R_+)}^2 \frac{\gamma'(u)}{\gamma(u)} du =c_h \|f\|_{L^2(\R_+)}^2\ .
\]
Similar arguments give, for the negative frequency part,
\[
\int_\R \|F(\cdot,u)\|_{L^2(\R_-)}^2 \frac{\gamma'(u)}{\gamma(u)} du =c_h \|f\|_{L^2(\R_-)}^2\ ,
\]
and putting both results together gives Equation~\eqref{fo:CQT_isometry}. Finally, Fatou's lemma gives the extension from $f\in C_c(\R)$ to $f\in L^2(\R)$.
\end{proof}
This result bears strong resemblance with known results on continuous wavelet transform~\cite{Grossmann1984decomposition}, in particular the admissibility condition. Notice however that the latter expresses a symmetry condition in the frequency domain with respect to frequency $\xi=-1$, while the corresponding wavelet admissibility condition expresses a similar symmetry with respect to the origin of frequencies. As a consequence, $h$ is necessarily complex-valued.

A closer connection can be made by introducing a function $\psi$ defined by
\begin{equation}
\label{fo:hr}
    \forall x\in\R,\quad \ \psi(x) = h(x)e^{2i\pi x}\ .
\end{equation}
where $\psi$ can be chosen real valued. Thus, the admissibility condition~\eqref{fo:CQT.admissibility} becomes
\begin{equation}
    0<c_\psi := \int_{\R_+}\frac{|\hat{\psi}(y)|^2}{y}dy = \int_{\R_-}\frac{|\hat{\psi}(y)|^2}{-y}dy<\infty\ ,
    \label{fo:CWT.admissibility}
\end{equation}
which is the usual admissibility condition for continuous wavelet transform~\cite{Grossmann1985transforms,Grossmann1986transforms}. We remind that the latter insures invertibility, a left inverse wavelet transform being given by the adjoint operator (up to the constant factor $c_\psi^{-1}$). 
The time-frequency atoms can then be written in terms of $\psi$ as
\begin{equation}
	\forall x,t,\u\in\R\ ,\quad \ h_{t,u}(x) := \sqrt{\gamma(u)}e^{2i\pi t}\psi(\gamma(u)x-t)\ ,
\end{equation}
which are closely related to wavelets as defined in~\cite{Grossmann1984decomposition}, with two mild modifications, namely the scale which is labeled by $\gamma(u)$, and a phase factor. These two changes do not modify strongly the classical wavelet transform.

\subsubsection{Transforms on $H^2(\R)$}
The constant-Q and wavelet transforms defined above turn out to be unsuitable for the construction we are about to describe. We found it more convenient to limit to functions whose Fourier transform vanishs for negative frequency. As in~\cite{Grossmann1984decomposition}, we introduce the real Hardy space
\begin{equation}
\label{es_real_hardy_space}
H^2(\R) = \left\{ f\in L^2(\R),\quad \hat f(\xi)=0\ \forall\xi\le 0\right\}\ .
\end{equation}
Let $\psi\in H^2(\R)$. Such a function $\psi$ is called analytic (or progressive) wavelet. The corresponding continuous wavelet transform~\cite{Grossmann1984decomposition} of a signal $f\in H^2(\R)$ is defined by
\begin{equation}
Wf(t,u) = \langle f,\psi_{t,u}\rangle = \frac{1}{\sqrt{\gamma(u)}}\,\int_{\R_+} \hat f(\xi)\overline{\hat\psi}\left(\frac{\xi}{\gamma(u)}\right) e^{2i\pi \xi t}\,d\xi\ ,\qquad u,t\in\R\ .
\end{equation} 
If the admissibility condition below is satisfied
\begin{equation}
\label{fo:CWT.admissibility.H2}
0<c_\psi := \int_{\R_+}\frac{|\hat{\psi}(y)|^2}{y}dy <\infty\ ,
\end{equation}
the corresponding transform satisfies the following isometry property
\begin{equation}
\left\|W f\right\|_{L^2(\R^2,d\mu)}^2 = c_\psi \|f\|^2_{H^2}\ ,
\end{equation}
and the measure is given by
\begin{equation}
d\mu(t,u) = \gamma'(u) du dt\ .
\end{equation}
\begin{remark}
The assumption $f\in H^2(\R)$ is not as irrelevant as it may appear. Indeed, in signal processing most signals are real-valued, so that their Fourier transform possess the Hermitean symmetry, \ie\, $\hat f(-\xi)=\overline{\hat f(\xi)}$. A real-valued signal $f\in L^2(\R)$ is then characterized by its orthogonal projection onto $H^2(\R)$, and can be reconstructed as the real part of the latter (up to a factor 2).
\end{remark}
\subsection{Definition of the frequency-focused transform}
\label{ssec:def.FFTFT}
We now introduce the frequency focus effect, generated by an associated frequency focus function $\sigma^\nu$. The role of the focus function is to modify the shape of the analysis waveforms, in a way that depends on some local behavior of the analyzed signal $f$. 

\paragraph{Assumptions}
Throughout this section, we make the following assumptions
\begin{enumerate}[i.]
\item 
$\psi\in H^2(\R)$ is an analytic wavelet function, therefore satisfying the admissibility condition~\eqref{fo:CWT.admissibility.H2}, and such that the quantity below (called frequency localization of $\psi$) is well-defined.
\label{FFTFT.assumption.psi}
\begin{equation}
\label{eq_wavelet_freq_loc}
\xi_0 := \frac{1}{\|\psi\|^2}\int_{\R_+} \xi |\hat\psi(\xi)|^2\,d\xi\ .    
\end{equation}
In addition, we assume that $|\hat\psi|^2$ is differentiable, and make the following technical assumptions:
\begin{itemize}
\item
$\hat\psi(\xi_0)\ne 0$
\item 
There exists $A_\psi>0$ such that for all $\xi\in\R_+$, 
\begin{equation}
\label{eq_hypothesis_psi}
\left(|\hat\psi|^2\right)'(\xi) \le \frac{A_\psi}{|\xi-\xi_0|}\ .
\end{equation}
\end{itemize}
%For technical reasons, we will make the following hypotheses \bruno{(toujours d'actualité ?)}\pierre{Oui normalement}
%\begin{equation}\label{eq_hypothesis_psi}
%\exists A_\psi,B_\psi >0 \text{ such that } \ \forall y\in\R, \  |y||\hat{\psi}(\xi_0+y)|\leqslant A_\psi \ \text{ and } \ \hat{\psi}'\in L^\infty
%\end{equation}
\item
$\gamma$ denotes a positive, strictly increasing $C^1$ diffeomorphism that maps $\R$ onto $\R_+^*$. 
\item
To every $f\in L^2(\R)$ is associated a focus function $\sigma^\nu_f$ of $f\in L^2(\R)$, assumed to be  continuous, larger than 1 and such that $\sigma_f^\nu -1$ goes to 0 at $\pm\infty$:
\begin{equation}
	\forall f\in L^2(\R), \quad \sigma^\nu_f \in C_{0}^+(\R) + 1\ .
\end{equation}
\end{enumerate}
%For every $u\in\R$, we introduce an affine function $\xi\to\beta_u(\xi)$, called squeezing function, of the form \bruno{(j'ai changé $\beta_\xi(u)$ en $\beta_u(\xi)$, ça me paraît plus naturel)}
%The latter, to be specified below, will control the frequency localization of time-frequency atoms. 

\paragraph{Time-frequency atoms and transform}
Wavelet ans constant-Q transforms use time-frequency atoms with constant relative bandwidth (\ie\, bandwidth divided by the frequency localization). The frequency-focused transform uses time-frequency atoms with prescribed frequency localization and bandwidth. This requires introducing an appropriate notion of frequency localization. Given a function $f\in H^2(\R)$, its frequency localization is defined by extending Equation~\eqref{eq_wavelet_freq_loc}:
$\frac{1}{\|f\|_{H^2}^2}\int_{\R_+} \xi |\hat f(\xi)|^2\,d\xi$
provided the integral is well-defined. 

\smallskip
The joint control of bandwidth and frequency localization is achieved by so-called squeezing functions $\xi\to\beta_u(\xi)$ defined as follows: for every $u\in\R$, we set
\begin{equation}
\label{eq_squeezing_fct}
\beta_u(\xi) = \frac{\sigma_f^\nu(u)}{\gamma(u)}\xi - \xi_1(u)\ , 
\end{equation}
for some shift parameter $\xi_1(u)>0$, to be specified below.

Given these parameters, we introduce frequency-focused atoms, defined by their Fourier transform
\begin{equation}
\widehat{\psi_{t,u,\sigma}}(\xi)=\frac{1}{\sqrt{\gamma(u)}}\,\hat\psi\left(\beta_u(\xi)\right)e^{-2i\pi\xi t}\ .
\end{equation}
A simple calculation shows that $\|\psi_{t,u,\sigma}\|^2 = \|\psi\|^2/\sigma(u)$.

The shift parameters $\xi_1(u)$ are fixed by imposing that the localization of $\left|\widehat{\psi_{t,u,\sigma}}(\xi)\right| $ equals $\gamma(u)\xi_0$, which yields
\begin{eqnarray*}
\gamma(u)\xi_0 &=& \frac{1}{\|\psi_{t,u,\sigma}\|^2}\frac{1}{\gamma(u)}
\int_{\R_+}\xi \left|\hat\psi\left(\frac{\sigma(u)}{\gamma(u)}\xi - \xi_1(u)\right)\right|^2\,d\xi\\
&=&\frac{1}{\gamma(u)\|\psi_{t,u,\sigma}\|^2}\int_{\R_+}\frac{\gamma(u)}{\sigma(u)}(\zeta+\xi_1(u))\left|\hat\psi(\zeta)\right|^2\frac{\gamma(u)}{\sigma(u)}\,d\zeta\\
&=&\frac{\gamma(u)}{\|\psi\|^2\sigma(u)}\left[\int_{\R_+}\zeta \left|\hat\psi(\zeta)\right|^2\,d\zeta + \xi_1(u) \|\hat\psi\|^2 \right]\\
&=&\frac{\gamma(u)}{\sigma(u)}\left[\xi_0+\xi_1(u)\right]\ ,
\end{eqnarray*}
therefore we obtain
\begin{equation}
\label{eq_xi_1}
\xi_1(u)=(\sigma(u)-1)\xi_0\ .
\end{equation}
Notice that when $u\to\pm\infty$, $\xi_1(u)\to 0$ and $\beta_u(\xi)\sim \xi/\gamma(u)$.

\medskip
The practical effect of such a squeezing is illustrated in Fig.~\ref{fi:freq_squeezing}, where a squeezing equal to 3 has been applied to three adjacent time-frequency atoms, whose bandwidth is therefore reduced while their amplitude is increased.

\begin{remark}
\begin{enumerate}
\item
Since $\xi_0>0$ and for all $u$ $\gamma(u)>0$ and $\sigma(u)\ge 1$, we have that $\beta_u(\xi)<0$ for all $\xi<0$; hence $\widehat{\psi_{t,u,\sigma}}(\xi)=0$ for all $\xi<0$.
\item
The frequency localization may actually be defined in several different ways. For example, assuming that $\hat\psi$ is a continuous function, the localization parameter $\xi_0$ may be defined as the mode of $|\hat\psi|$, by setting $\xi_0 = \mathop{\text{argmax}}_{\xi\in\R_+^*} |\hat\psi(\xi)|$.
In this case, using the same localization measure for $\widehat{\psi_{t,u,\sigma}}$, imposing that $\mathop{\text{argmax}}_{\xi\in\R_+^*} |\widehat{\psi_{t,u,\sigma}}| = \xi_0\gamma(u)$ is equivalent to $\beta_u(\gamma(u)\xi_0)=\xi_0$, which yields the same expression~\eqref{eq_xi_1} for the frequency shifts $\xi_1(u)$.
\end{enumerate}
\end{remark}


\begin{figure}[h!]
    \centering
    \includegraphics[scale=.5]{freq_squeezing}
    \caption{Frequency squeezing: Fourier transforms of constant-Q atoms (top), and the same atoms out of which three have been squeezed (bottom) by a factor 3.}
    \label{fi:freq_squeezing}
\end{figure}

Given these notations, the frequency-focused transform $M^\nu$ can be defined for $f\in C_c$ by
\begin{equation}
    \forall t,u\in\R,\quad \ M^\nu f(t,u) := \langle f , \psi_{t,u,\sigma_f^\nu}\rangle_{L^2}\ ,
\end{equation}
Plancherel's formula gives the following form
\begin{equation}
M^\nu f(t,u) = \langle \hat f,\widehat{\psi_{t,u,\sigma_f^\nu}}\rangle = \frac{1}{\sqrt{\gamma(u)}}\,\int_{\R_+} \hat f(\xi)\overline{\hat\psi}\left(\frac{\sigma_f^\nu(u)}{\gamma(u)}\xi - (\sigma(u)-1)\xi_0\right) e^{2i\pi \xi t}\,d\xi
\end{equation}


\begin{remark}
Notice that because of our choice of normalization, the time-frequency atoms do not have constant norm any more. Retaining constant norm would impose to normalize them by $\sqrt{\sigma_f}$ instead of $\sigma_f$, but this would in turn lead to multiply the measure $\mu$ by a factor that depends explicitly on $f$, which we want to avoid.
\end{remark}

\paragraph{Motivations and examples for the frequency focus function}
Examples of time focus functions were given in the corresponding paragraph in the previous Section. The rationale for frequency focus functions should follow similar objectives, we won't discuss them here, and refer to Section~\ref{se:numerical.experiments} devoted to numerical experiments.

\bigskip

The main purpose of the next sections is to establish that $M^\nu$ can be extended to a well-defined map from $H^2(\R,dx)\rightarrow L^2(\R^2,d\mu)$ satisfying a norm control similar to a frame bound control.  The pointwise definition of $M^\nu f(t,\u)$ on $L^2$ is still guaranteed by the fact that $h\in L^2$. Regarding the definition from $H^2$ into $L^2(\R^2,d\mu)$ we can raise that $M^\nu$ is well-defined from $C_c$ into $L^2(\R^2,d\mu)$ and then extend the control by the use of the Fatou's lemma.

Sometimes, if there are no ambiguities, we will write $\psi_{t,\u,f}$ instead of $\psi_{t,\u,\sigma_f}$ and $\psi_{t,\u,n}$ instead of $\psi_{t,\u,\sigma_{f_n}}$ for a certain sequence $(f_n)_n$.


\begin{comment}
\subsection{Example and \review{motivation} of the frequency focus function}\label{ssec_frequency_focus_function}
Defining and studying specific frequency focus functions is not a main goal of this paper. We can however provide a toy example that could be of interest in some specific situations.
\begin{example}[Resolving close frequencies]
\label{ex:frequency_focus}
	Imagine a function $f$ of interest contains several harmonic components (i.e. sine waves) with more or less close frequencies. When consecutive frequencies are too close, a wideband window will not be able to separate them, which would suggest increasing the frequency precision in the neighborhood of these frequencies. A prototypical choice for frequency focus function in such a context could be
	\begin{equation}
 \label{fo:ex.freq.foc}
		\sigma_f^\nu(\u) = 1 + \frac{\|V_gf(\cdot,\u)\|_{L^1(\R)}}{\|V_gf(\cdot,\u)\|_{L^2(\R)}}\,\varphi(\u)\ ,
	\end{equation}
	where $\varphi$ is compactly supported in some interval $[-\u_0,\u_0]$ for some cutoff frequency $\u_0$ ($\varphi$ can be a multiple of the indicator function of that interval, or the hearing threshold in psycho-physical applications), and $V_g$ is the ordinary STFT, with Gaussian window~\cite{grochenig2013foundations}.\\
	%
	The ratio of $L^1$ to $L^2$ norms is known as a a measure of spread, or diversity (it takes small values when the corresponding function is more concentrated, and is closely related to a specific Renyi entropy, see~\cite{Ricaud2013refined}). When $f$ contains several sine waves in the neighborhood of $\u$, which are not separated by the \review{Gaussian} window, some beats appear in the function $t\mapsto V_gf(t,\u)$, which yields smaller $L^1/L^2$ norm values, and therefore small values for $\sigma_f^\nu(\u)$.

\smallskip
Illustrations with a suitably post-processed version of this choice, on a toy signal example, is provided in Appendix~\ref{sec:app.freq.foc.example}.
%
\end{example}
\end{comment}
\subsection{Kernel and norm relationship}

To prove the main result we first derive a norm relationship involving a certain non-negative valued kernel $K_\sigma$, so that the study of the norm of $M^\nu f$ will be determined by the norm $\|K_\sigma\|_\infty$.
\begin{prop}
	Let $\psi$ satisfying assumptions~\ref{FFTFT.assumption.psi}. in Section~{\ref{ssec:def.FFTFT}}. Then we have for every $f\in H^2(\R)$
	\begin{equation}
		\|M^\nu f\|_{L^2(d\mu)}^2 = \int_{\R_+} |\hat{f}(\xi)|^2 K_{\sigma_f^\nu}(\xi)d\xi\ ,
	\end{equation}
	where the kernel $K_{\sigma_f^\nu}(\xi)$ is defined by 
\begin{equation}\label{kernel_frequency_case}
		K_{\sigma_f^\nu}(\xi) := \int_{\R_+} \left|\hat{\psi}\left( \beta_{\gamma^{-1}(y)}(\xi) \right)\right|^2\frac{dy}{y}\ .
\end{equation}
\end{prop}
%\begin{remark}
%\review{Since $\sigma\circ\gamma^{-1}(|\frac{\xi}{y}|)=1$ for $y$ big enough or close enough to $0$, thanks to the admissibility condition on $\psi$  (see Equation~\eqref{fo:CWT.admissibility}) the convergence of $K_\sigma$ is guaranteed. \bruno{à revoir ?}}
%\end{remark}
\begin{proof}
We first introduce the auxiliary function $F(\xi,u) = \hat f(\xi) \overline{\hat\psi}(\beta_u(\xi))$, and notice that $F(\cdot,u)\in L^1(\R_+)$ for all $u\in\R$. Then compute
\begin{eqnarray*}
\|M^\nu f\|_{L^2(\R^2,d\mu)}^2&=& \int_{\R^2\times\R_+^2} F(\xi,u)\overline{F}(\xi',u) e^{2i\pi (\xi-\xi')t}\frac{\gamma'(u)}{\gamma(u)}d\xi d\xi' dudt\\
&=&\int_{\R^2} \left|(\mathcal{F}_1^{-1} F)(t,u)\right|^2 \frac{\gamma'(u)}{\gamma(u)} dudt\\
&=& \int_{\R^2} \left| F(\xi,u)\right|^2 \frac{\gamma'(u)}{\gamma(u)} dud\xi\\
&=& \int_{\R^2} \left|\hat f(\xi)\right|^2\,\left|\hat\psi(\beta_u(\xi))\right|^2 \frac{\gamma'(u)}{\gamma(u)} dud\xi\\
&=&\int_{\R_+^2} \left|\hat f(\xi)\right|^2\,\left|\hat{\psi}\left( \beta_{\gamma^{-1}(y)}(\xi) \right)\right|^2\frac{dy}{y}d\xi\ ,
\end{eqnarray*}
where we have denoted by $\mathcal{F}_1^{-1}F$ the inverse Fourier transform of $F$ with respect to its first variable, and then used the corresponding Plancherel formula. The argument above involve the use of Fubini's theorem which is justified by the fact that the integral with respect to $y$ is convergent, this fact is proved in Theorem~\ref{theorem_norm_control_frequency_focus}.
\end{proof}

\begin{comment}
\begin{proof}
Introducing $\sigma(\u)$ in the proof of Proposition~\ref{prop_norm_relation_operator_Q_transform} gives us 
	\begin{align*}
		\|M^\nu f\|^2_{L^2} & = \int_{\review{\R}} \bar{\hat{f}}(\xi)\hat{f}(\xi')\frac{1}{\gamma(\u)}\hat{\psi}\left(\tfrac{\xi'}{\sigma_f^\nu(u)\gamma(u)}\right)\bar{\hat{\psi}}\left(\tfrac{\xi}{\sigma_f^\nu(u)\gamma(u)}\right)e^{-2i\pi\gamma(\u)^{-1}(\xi'-\xi)t}d\xi d\xi'd\mu(\u,t) \\
		& = \int_{\review{\R}}|\hat{f}(\xi)|^2 \frac{\gamma'(\u)}{\gamma(\u)^2} \left|\hat{\psi}\left(\tfrac{\xi}{\sigma_f^\nu(u)\gamma(u)}\right)\right|^2 d\xi d\u \\
		& = (1) + (2)\ ,
	\end{align*}
	where
	\begin{align*}
		(1) &:= \int_0^{+\infty}|\hat{f}(\xi)|^2\int_{-\infty}^{+\infty}\frac{\gamma'(\u)}{\gamma(\u)^2} \left|\hat{\psi}\left(\tfrac{\xi}{\sigma_f^\nu(u)\gamma(u)}\right)\right|^2 d\u d\xi, \\
		(2) & := \int_{-\infty}^0|\hat{f}(\xi)|^2\int_{-\infty}^{+\infty}\frac{\gamma'(\u)}{\gamma(\u)^2} \left|\hat{\psi}\left(\tfrac{\xi}{\sigma_f^\nu(u)\gamma(u)}\right)\right|^2 d\u d\xi\ .
	\end{align*}
	Using the same change of variable $y= \tfrac{\xi}{\gamma(u)}$ in both cases and taking care of the signs we get (1) and in (2) the sign of $\xi/y$ is still non-negative as $\xi\leqslant 0$ and $y\leqslant 0$ so we obtain
	\begin{align*}
		(1) &=  \int_0^{+\infty}|\hat{f}(\xi)|^2\int_{0}^{+\infty} \left|\hat{\psi}\left( \frac{y}{\sigma_f^\nu\circ\gamma^{-1}(\frac{\xi}{y})} \right)\right|^2\frac{dy}{y} d\xi , \\
		(2) &= -\int_{-\infty}^0|\hat{f}(\xi)|^2 \int_{-\infty}^{0} \left|\hat{\psi}\left( \frac{y}{\sigma_f^\nu\circ\gamma^{-1}(\frac{\xi}{y})} \right)\right|^2\frac{dy}{y} d\xi \\
		&= \int_{-\infty}^0|\hat{f}(\xi)|^2 \int_{-\infty}^{0} \left|\hat{\psi}\left( \frac{y}{\sigma_f^\nu\circ\gamma^{-1}(|\frac{\xi}{y}|)} \right)\right|^2\frac{dy}{|y|} d\xi.
	\end{align*}
	Using a symmetry argument similar to the proof of Proposition~\ref{prop_norm_relation_operator_Q_transform}, and summing both terms we obtain the expected formula.
\end{proof}
\end{comment}
A simple change of variable gives the following alternative expression for the kernel:
\begin{corollary}\label{cor_second_expression_kernel}
%    With a change of variable we have the other way to write the kernel,
The kernel $K_\sigma$ may be written as
    \begin{equation}
        \forall \xi>0 \ ,\quad \ K_\sigma(\xi)= \int_{\R_+}\left|\hat{\psi}\left(\sigma\circ\gamma^{-1}\left(\tfrac{\xi}{y}\right)y-\xi_0\left(\sigma\circ\gamma^{-1}\left(\tfrac{\xi}{y}\right)-1\right)\right)\right|^2\tfrac{dy}{y}\ .
    \end{equation}
\end{corollary}

	\begin{remark}
		If $\sigma$ is fixed and independent of $f$, the time-frequency atoms $h_{t,\omega,\sigma(\omega)}$ form a \textit{continuous frame} of $L^2(\R)$ in the terminology of~\cite{Ali2000coherent,Dahlke2008generalized}. We will see that in the general case, the assumptions made on $\sigma$ allow one to stay in a tractable situation.
\end{remark} 
From now on, the goal is to obtain upper and lower bounds for the kernel $K_\sigma$.

\subsection{Main result : the $L^2$ norm control}\label{ssection_norm_control_frequency_focus}

The main result of this section is the following Theorem~\ref{theorem_norm_control_frequency_focus} which is a generalization to non-linear transform with adaptive window of the classical frame control.
\begin{theorem}\label{theorem_norm_control_frequency_focus}
Let $\psi$ satisfying assumptions~\ref{FFTFT.assumption.psi}. in Section~{\ref{ssec:def.FFTFT}}.
	Let $f\in H^2(\R)$,
 %$\psi$ satisfying hypothesis~\eqref{eq_hypothesis_psi}
 and let $\sigma^\nu_f\in C_0^+(\R)+1$ denote the associated frequency-focus function.
 %, lets $0<a<b<+\infty$ such that
    Then we have
	\begin{equation}
		d_\psi\|f\|_{L^2(\R)}^2 \leqslant \|M^\nu f \|_{L^2(\R^2,d\mu)}^2 \leqslant C_{\sigma_f^\nu}\|f\|_{L^2(\R)}^2\ ,
	\end{equation}
	where $d_\psi>0$ depends only on the wavelet $\psi$, and $C_{\sigma_f^\nu}$ is given by
	\begin{equation}
		C_{\sigma_f^\nu} := c_\psi + A_\psi\ \int_{\R_+}\left(\sigma_f^\nu(\gamma^{-1}(y))-1\right)\frac{dy}{y} <+\infty,
	\end{equation}
%	\begin{equation}
%		c_{a,b}^2 := \frac{|\hat{\psi}(\xi_0)|^2}{2}\ln\left(\frac{b}{a}\right).
%	\end{equation}
\end{theorem}
The Theorem is proven in Propositions~\ref{prop_upper_bound_control} and~\ref{prop_lower_bound_kernel_basic}.
\begin{remark}
\label{rem:FFTFT.lower.bound}
%The fact that $d_\psi$ does not depend on the focus function $\sigma_f^\nu$, and therefore $f$,  will insure the injectivity of the transform.
As we shall see in the proof below, the assumptions on $\psi$ insure the existence of an interval $(a,b)$ containing $\xi_0$ such that $|\hat\psi(\xi)|^2 \ge |\hat\psi(\xi_0)|^2/2$ for all $\xi\in (a,b)$, which yields the lower bound
\[
d_\psi = \frac{|\hat{\psi}(\xi_0)|^2}{2}\ln\left(\frac{b}{a}\right)\ .
\]
\end{remark}
%\begin{remark}
%	We can raise that if $\sigma$ is fixed (\ie\, independent of $f$) then we obtain the classical bound control from linear time frequency operators since $c_{a,b}$ and $C_{a,b}$ only depend on $\sigma$ and not directly on $f$. More precisely they only depend on the support of $\sigma$.
%\end{remark}

\begin{comment}
\begin{proof}
	Even though $\sigma_f$ is supposed to be \review{piecewise} continuous, the idea of the proof is to control the sup norm of the kernel $K_\kappa$ when $\kappa \in \mathcal{E}_c^+ +1$ with $\mathcal{E}$ the set of step functions and $\kappa$ close enough to $\sigma$. Everything will be defined correctly in the next section. The motivation is the fact that the kernel $K_\kappa$ is easy to estimate when $\kappa$ is a step function. Thus, by using the control for step functions and proving a continuity result of $K_\kappa$ for functions $\kappa\rightarrow \sigma$ we will be able to prove our control for continuous (or piecewise continuous) focus functions $\sigma_f$.
\end{proof}
\end{comment}

%\subsection{Proof of Theorem \ref{theorem_norm_control_frequency_focus}}
%\label{ssection_rectangular_function_frequency_focus}

\subsubsection{Upper bound control}

We have the following upper bound control.

\begin{prop}\label{prop_upper_bound_control}
    Let $\sigma\in C_0^+(\R)+1$ be a focus function.
    \begin{align}
        \|K_{\sigma}\|_\infty \leqslant c_\psi + A_\psi \int_{\R^+}\left(\sigma(\gamma^{-1}(y))-1\right)\frac{dy}{y}\ ,
    \end{align}
    where $A_\psi$ is given in Equation~\eqref{eq_hypothesis_psi}.
\end{prop}
\begin{proof}
    Let $\xi\geqslant0$ and write $\vphi(\xi)=|\hat{\psi}(\xi)|^2$ for simplicity. We will use the following expression for the kernel, which results from a change of variable,
    \begin{equation*}
        K_\sigma(\xi) = \int_\R \vphi(\beta_u(\xi))\frac{\gamma'(u)}{\gamma(u)}du.
    \end{equation*}
    Hence, by the mean value theorem, we can write
    \begin{align*}
        K_\sigma(\xi) - c_\psi & = \int_\R\left[ \vphi(\beta_u(\xi)) - \vphi\left(\tfrac{\xi}{\gamma(u)}\right) \right]\tfrac{\gamma'(u)}{\gamma(u)} du \\
        & = \int_\R \left( \beta_u(\xi)-\tfrac{\xi}{\gamma(u)} \right) \vphi'(\zeta_{u,\xi})\tfrac{\gamma'(u)}{\gamma(u)}du\ ,
    \end{align*}
    with
    \begin{equation*}
        \zeta_{u,\xi} = \xi_0 +\left( \tfrac{\xi}{\gamma(u)}-\xi_0 \right)(1+\theta_{u,\xi}(\sigma(u)-1)), \quad \theta_{u,\xi}\in (0,1).
    \end{equation*}
    Since $\beta_u(\xi)-\tfrac{\xi}{\gamma(u)}=(\sigma(u)-1)\left( \tfrac{\xi}{\gamma(u)}-\xi_0 \right)$ and using hypothesis \eqref{eq_hypothesis_psi} we obtain
    \begin{align*}
        |K_\sigma(\xi) - c_\psi| & \leqslant \int_\R (\sigma(u)-1)\left|\left( \tfrac{\xi}{\gamma(u)}-\xi_0 \right)\vphi'\left( \xi_0 +\left( \tfrac{\xi}{\gamma(u)}-\xi_0 \right)(1+\theta_{u,\xi}(\sigma(u)-1)) \right) \right|\tfrac{\gamma'(u)}{\gamma(u)} du \\
        & \leqslant \int_\R (\sigma(u)-1)\frac{A_\psi}{1+\theta_{u,\xi}(\sigma(u)-1)}\tfrac{\gamma'(u)}{\gamma(u)}du \\
        & \leqslant A_\psi \int_\R (\sigma(u)-1)\tfrac{\gamma'(u)}{\gamma(u)}du \\
        & =A_\psi \int_{\R^+}\left(\sigma(\gamma^{-1}(y))-1\right)\frac{dy}{y}\ ,
    \end{align*}
    which achieves the proof of the Proposition.
\end{proof}

\subsubsection{Lower bound control}

The following result guarantees the existence of a positive lower bound that only depends on the wavelet $\psi$. Under the hypothesis~\ref{FFTFT.assumption.psi}. in Section~\ref{ssec:def.FFTFT}, we have the existence of $a<b\in\R_+^*$ such that
\begin{equation}
        \forall y\in (a,b),\ |\hat{\psi}(y)|^2 \geqslant \frac{|\hat{\psi}(\xi_0)|^2}2\ .
\end{equation}
We can then prove
\begin{prop}\label{prop_lower_bound_kernel_basic}
Let $\sigma\in C_0^+(\R)+1$ be a focus function. 
Then
    \begin{equation}
        \forall \xi>0, \ K_{\sigma}(\xi)\geqslant \frac{|\hat{\psi}(\xi_0)|^2}{2}\ln\left(\frac{b}{a}\right).
    \end{equation}
\end{prop}

\begin{proof}
    Let us fix $\xi>0$. Using the notation 
    \[
    \alpha_\xi(y)=\beta_{\gamma^{-1}(y)}(\xi) = \left(\sigma\circ\gamma^{-1}\right)(y)\tfrac{\xi}{y}-\xi_0\left(\left(\sigma\circ\gamma^{-1}\right)(y)-1\right)\ ,
    \]
    we can write
    \begin{equation*}
        K_\sigma(\xi)= \int_{\R_+} |\hat{\psi}(\alpha_\xi(y))|^2\frac{dy}{y}\ .
    \end{equation*}
    Since $\alpha_\xi$ is continuous, we have
    \begin{align*}
        \alpha_\xi((\xi/b,\xi/a))&\supset \left(\alpha_\xi(\xi/b),\alpha_\xi(\xi/a)\right) \\
        & = \left((a-\xi_0)\left(\sigma\circ\gamma^{-1}\right)(1/a)+\xi_0 , (b-\xi_0)\left(\sigma\circ\gamma^{-1}\right)(1/b)+\xi_0 \right)\\
        & \supset (a,b)\ .
    \end{align*}
    Indeed, since $a-\xi_0<0$ and $b-\xi_0>0$, together with the fact that $\sigma(u)\ge 1$ for all $u$, we have 
    \begin{align*}
        (a-\xi_0)\left(\sigma\circ\gamma^{-1}\right)(1/a) &\leqslant a-\xi_0 \ , \\
        (b-\xi_0)\left(\sigma\circ\gamma^{-1}\right)(1/b) &\geqslant b-\xi_0 \ .
    \end{align*}
    Hence
    \begin{align*}
        \int_{\R^+} |\hat{\psi}(\alpha_\xi(y))|^2\frac{dy}{y} &\geqslant \frac{|\hat{\psi}(\xi_0)|^2}{2}\int_{\alpha_\xi^{-1}((a,b))}\frac{dy}{y} \\
        & \geqslant \frac{|\hat{\psi}(\xi_0)|^2}{2} \int_{\xi/b}^{\xi/a}\frac{dy}{y} \\
        & \geqslant \frac{|\hat{\psi}(\xi_0)|^2}{2}\ln\left(\frac{b}{a}\right)\ ,
    \end{align*}
    which proves the proposition, and yields the expression of the bound given in Remark~\ref{rem:FFTFT.lower.bound}.
\end{proof}



